\chapter{Transcranial Magnetic Stimulation}

\section*{Overview}
Transcranial magnetic stimulation (TMS) is a noninvasive technique that uses magnetic fields to stimulate neurons in targeted regions of the brain. It is used both as a diagnostic tool to assess cortical excitability and as a therapeutic treatment, most commonly for depression.

\section*{Alternative Names}
TMS, repetitive transcranial magnetic stimulation (rTMS), theta burst stimulation (TBS), motor evoked potential testing

\section*{Principle}
A rapidly changing magnetic field passes through the skull and induces an electrical current in underlying cortical neurons, depolarizing them. Repetitive stimulation (rTMS) at low frequencies tends to suppress cortical activity, whereas high frequencies increase it. Repeated sessions of stimulation to the dorsolateral prefrontal cortex are thought to modulate dysfunctional circuits underlying mood disorders.

\section*{History}
Barker and colleagues developed modern TMS at the University of Sheffield in 1985. In 2008, rTMS received FDA approval for the treatment of medication-resistant major depressive disorder.

\section*{Why This Test or Procedure Is Ordered}
Therapeutically ordered for major depressive disorder that has not responded adequately to one or more antidepressant trials, and for several other conditions including obsessive-compulsive disorder and migraine. Diagnostically, it is used to map motor cortex function and to test corticospinal tract integrity.

\section*{Is This a Primary or Secondary Test or Procedure}
A secondary treatment used after inadequate response to pharmacotherapy; a diagnostic adjunct to routine neurologic evaluation.

\section*{How This Test or Procedure Is Performed}
The patient is awake and seated in a chair while a figure-of-eight coil is placed against the scalp. The resting motor threshold is first determined by observing thumb or hand twitch. Treatment delivers trains of repeated pulses to the target region in sessions lasting 20-40 minutes, typically five days per week for four to six weeks.

\section*{What Preparation Is Needed}
Screening for metallic implants and history of seizures, and removal of metal objects, hearing protection, and jewelry. No sedation is required, and the patient may eat and take regular medications.

\section*{Contraindications and Precautions}
Contraindicated in patients with ferromagnetic metal in the head (aneurysm clips, cochlear implants) or implanted devices affected by magnetic fields, such as vagus nerve or deep brain stimulators. Precautions include seizure history, epilepsy, and pregnancy.

\section*{Risks}
The most common side effects are mild headache, scalp discomfort, and facial twitching that resolve rapidly. The most serious risk is seizure, which is rare, along with transient hearing changes or mania.

\section*{What Happens After the Test or Procedure}
Patients may immediately resume normal activities and driving. Response is assessed over the treatment course with symptom rating scales, typically evident after two to four weeks.

\section*{Understanding Your Results}
Therapeutic response is measured clinically rather than by imaging or laboratory values, using standardized depression scores. Maintenance sessions may be scheduled for patients who respond, to sustain the benefit.

\section*{Normal Results}
Reduction in depression severity scores, such as a 50\% or greater drop on the Hamilton Depression Rating Scale. Recovery of motor evoked potentials with normal latencies reflects intact corticospinal conduction.

\section*{Follow-up Tests and Procedures}
Serial mood rating scales during and after the treatment course to document response. Responders may require tapering maintenance rTMS or continued antidepressant therapy to prevent relapse.

\section*{References}
\begin{enumerate}
  \item MedlinePlus. Transcranial magnetic stimulation. U.S. National Library of Medicine; 2025.
  \item Current Medical Diagnosis and Treatment. McGraw-Hill; 2025.
\end{enumerate}

\clearpage
